%----------------------------------------
% IT IS RECOMMENDED TO USE AUTOLATEX FOR
% COMPILING THIS DOCUMENT.
% http://www.arakhne.org/autolatex
%----------------------------------------

\documentclass[english,nodocumentinfo,nopubpage]{utbmciadreport}

% The TeX code is entering with UTF8
% character encoding (Linux and MacOS standards)
\usepackage[utf8]{inputenc}
\usepackage{fancyhdr}
\usepackage{../common/sarl-colorized-listing}


\graphicspath{{./}}

\declaredocument{Introduction to SARL}{Tutorial Session \#6}{UTBM-INFO-AI51-TU6}

\addauthorvalidator*[St\'ephane Galland]{St{\'e}phane}{Galland}{Teacher}

\updateversion{1.0}{\makedate{23}{03}{2023}}{First release on Github}{\upmpublic}

\begin{document}

\section{Goal of this Tutorial Session}

The goal of this tutorial session is to introduce the basics concepts of agent-oriented programming with SARL.
The general goal of the tutorial is to remind several key points related to the coding with SARL.

\section{Work to be Done during the Tutorial Session}

The following sections describe the work to be done during this tutorial session. It is recommended to work 30 minutes on each exercice.
To do these exercice, it is recommended to read the online documentation of SARL for finding the best functions or capacities to be used.

\subsection{Exercice 1 - Ping-Pong Communication Protocol}

In this exercice, the goal is to implement a basic communication protocol between two agents: the Ping-Pong.
The principle of this protocol is the following:
\begin{enumerate}[1.]
\item The first agent is creating a \code{Ping} event and sent it to the second agent;
\item When the second agent is receiving a \code{Ping}\ event, the agent has to replies to the emitted of the \code{Ping}\ event with a specific \code{Pong}\ event.
\item When the first agent is receiving a \code{Pong}\ event, the agent has to replies to the emitted of the \code{Pong}\ event with a specific \code{Ping}\ event.
\end{enumerate}

The steps to follow for obtaining the answers are:
\begin{enumerate}
\item Define the event with the name \code{Pong}.
\item Define the event with the name \code{Ping}.
\item Define the agent with the name \code{Agent1}, that is firing a \code{Ping} event to any other agent.
\item Define the agent with the name \code{Agent2}, that is firing a \code{Pong}\ event when it receives a \code{Ping}\ event. The receiver of the \code{Pong}\ event must be the emitter of the just received \code{Ping}\ event.
\item Update the code of \code{Agent1} for firing a new \code{Ping}\ event when it receives a \code{Pong} event. The receiver of the \code{Ping}\ event must be the emitter of the just received \code{Pong} event.
\end{enumerate}

You have to start the two agents to run the Ping-Pong properly. How to start two agent?

\subsection{Exercice 2 - Extended Ping-Pong Communication Protocol}

In the previous exercice, the spawning of the agents, and the fact that the first agent is emitting the initial \code{Ping}\ event when it is started, it may cause issues.

\begin{enumerate}
\item What is this possible issue in the previous excercice?
\item Solve the issue by ensuring in the structure and the behavior of your code that the initial occurrence of the \code{Ping}\ event is fired only when the second agent is alive.
\end{enumerate}

\subsection{Exercice 3 - Contract Net Protocol}

In this excercise, you must define a multiagent system that is running a Contract Net Protocol:
\begin{enumerate}
\item A customer wants to by a product.
\item The customer sends a request to a broker for finder the best provider.
\item The broker asks to a couple of provider they offers for the product.
\item The broker selects the best offer.
\item The broker notifies the selected provider about its acceptance with the ID of the customer.
\item The broker notifies the not selected providers about their rejections.
\item The broker notifies the customer with the ID of the provider.
\end{enumerate}

You should:
\begin{enumerate}[a)]
\item Define the capacity of the provider.
\item Define the capacities of the broker.
\item Write the provider code.
\item Write the broker code.
\item Write the customer code.
\end{enumerate}

\subsection{Exercice 4 - Basic Simulation Loop with SARL}

In this excercise, you will learn how to implement a basic simulation loop.
The multiagent system is composed by three agents:
\begin{itemize}
\item An agent \code{AgentEnvironment}\ that is representing the agent environment (see Course Slides).
\item Two agents of type \code{ApplicationAgent}\ that are representing the agents that are doing actions in the agent environment.
\end{itemize}

For implementing the basic simulation loop, you have to follow the steps below:
\begin{enumerate}
\item Define an event \code{Perception}\ that represents a data that must be used by the \code{ApplicationAgent} to decide to do something.
\item Define an event \code{Action}\ that represents an action that an \code{ApplicationAgent}\ would do in the \code{AgentEnvironment}. This event must contain a string of characters that is description the action to be done.
\item Define the type \code{AgentEnvironment}\ that have the following behavior:
	\begin{itemize}
	\item Define a list of \code{UUID}, named \code{bodies}. This list contains the UUID (identifier of the agent) of all the agents of type \code{ApplicationAgent}\ that are spawned in the system.
	\item Define the function \code{runStartOfLoop}\ that is looping on the \code{bodies}\ list and send an event of type \code{Perception}\ to each agent with the UUID (the identifier of the agent).
	\end{itemize}
\item Define the type \code{ApplicationAgent}\ that have the following behavior:
	\begin{itemize}
	\item When the agent receives an event of type \code{Perception}, the agent:
		\begin{enumerate}[a)]
		\item Create an instance of the event \code{Action}\ with the action decided by the agent described in the attribute of the event.
		\item Fire the instance of the event \code{Action}\ to the agent of type \code{AgentEnvironment}.
		\end{enumerate}
	\end{itemize}
\item Update the type \code{AgentEnvironment}\ with the following behavior:
	\begin{itemize}
	\item Create an event handler that is corresponding to the receiving of the event \code{Action}.
	\item The previous event handler must print out the action decided by an \code{ApplicationAgent}\ that is received by the \code{AgentEnvironment}.
	\item When the \code{AgentEnvironment}\ has received the same number of action events as the number of UUID in the \code{bodies}\ attribute, then the actions of the agents must be applied in the environment and the simulation loop must restart. For restarting the simulation loop, call the function \code{runStartOfLoop}.
	\end{itemize}
\item Update the type \code{AgentEnvironment}\ with the following behavior:
	\begin{itemize}
	\item At startup, create a given number of \code{ApplicationAgent} and save the UUID of the spawned agents in the \code{bodies}\ attributes.
	\item Start the simulation loop \Emph{asynchronously}.
	\end{itemize}
\end{enumerate}	

\end{document}
